\documentclass[letterpaper,12pt]{article}

% @@@@@@@@@@@@@@@@@@@@@@@@@@@@@@@@@@@@@@@@@@@@@@@@@@@@@@@@@@@@>
% VALORES A MODIFICAR POR USTED:
% @@@@@@@@@@@@@@@@@@@@@@@@@@@@@@@@@@@@@@@@@@@@@@@@@@@@@@@@@@@@>
\newcommand{\titulo}{Título de la memoria}
\newcommand{\ciudad}{Ciudad} % e.g. Valparaíso
% TODO: Consultar el formato de los nombres:
\newcommand{\nombrealumno}{Nombre del alumno}
\newcommand{\nombreprofesor}{Nombre del profesor guía}
\newcommand{\nombrecorreferente}{Nombre del correferente}
% Mes y año del examen
\newcommand{\mesexamen}{Diciembre}
\newcommand{\anioexamen}{20XX}
% Dedicatoria y agradecimientos
\newcommand{\dedicatoria}{
Considerando lo importancia de este trabajo para los alumnos, este apartado es para que el autor entregue palabras personales para dedicar este documento. La extensión puede ser de máximo una hoja y se deben mantener este formato, tipo y tamaño de letra.
}
\newcommand{\agradecimientos}{
Considerando la importancia de este trabajo para los alumnos, este apartado se podrá incluir en el caso de que el autor desee agradecer a las personas que facilitaron alguna ayuda relevante en su trabajo para la realización de este documento. La extensión puede ser de máximo una hoja y se deben mantener este formato, tipo y tamaño de letra.
}
\newcommand{\resumen}{
El resumen y las palabras clave no deben superar la mitad de la página, donde debe precisarse brevemente: 1) lo que el autor ha hecho, 2) cómo lo hizo (sólo si es importante detallarlo), 3) los resultados principales, 4) la relevancia de los resultados. El resumen es una representación abreviada, pero comprensiva de la memoria y debe informar sobre el objetivo, la metodología y los resultados del trabajo realizado.
}
\newcommand{\resumeningles}{
Corresponde a la traducción al idioma inglés del Resumen anterior. Sujeto a la misma regla de extensión del Resumen.
}
\newcommand{\palabrasclave}{
Cinco es el máximo de palabras clave para describir los temas tratados en la memoria, ponerlas separadas por punto y comas.
}
\newcommand{\palabrasclaveingles}{
Corresponde a la traducción al idioma inglés de Palabras Clave anteriores.
}
% @@@@@@@@@@@@@@@@@@@@@@@@@@@@@@@@@@@@@@@@@@@@@@@@@@@@@@@@@@@@>

% Idioma y fuentes
\usepackage[spanish]{babel}
\usepackage[T1]{fontenc}
\usepackage{fontspec}

% Paquete para definir cualquier tamaño de font
\usepackage{anyfontsize}

% NOTE: La font principal, debería ser Calibri, pero hay que pagar por ella al menos que se esté usando Windows. En su defecto se usa la métricamente compatible Carlito. Si (de alguna manera) tiene instalada Calibri en su sistema, puede colocarla aquí:
\newcommand{\fuente}{Carlito}

\setmainfont{\fuente}

% Tamaño de la página y márgenes
\usepackage[letterpaper,top=2.5cm,bottom=3cm,left=3cm,right=3cm,marginparwidth=1.75cm]{geometry}

% Determinar interlineado:
\renewcommand{\baselinestretch}{1.0}

% Eliminar sangrías:
\setlength{\parindent}{0cm}

% Paquete para definir los formatos de los títulos
\usepackage[explicit]{titlesec}

\titleformat{name=\section}[block]{\fontsize{16}{24}\selectfont\bfseries}{}{0pt}{#1}
\titleformat{name=\section,numberless}[block]{\fontsize{16}{24}\selectfont\bfseries}{}{0pt}{#1}
\titlespacing*{name=\section}{0pt}{0pt}{0.5cm}
\titlespacing*{name=\section,numberless}{0pt}{0pt}{0.5cm}

% Separación entre parrafos
\setlength{\parskip}{0.4cm}

% Paquetes de utilidad general
\usepackage{amsmath}
\usepackage{graphicx}
\usepackage[colorlinks=true, allcolors=blue]{hyperref}

% Formato de las tablas de contenido
\usepackage[tocflat]{tocstyle}
\usetocstyle{allwithdot}

% Para obtener el número de la última página
\usepackage{lastpage}

% Header y footer
\usepackage{fancyhdr}
\fancypagestyle{portada}{
    \lhead{}
    \chead{}
    \rhead{}
    \lfoot{}
    \cfoot{\fontsize{10}{12}\selectfont \thepage}
    \rfoot{}
    \renewcommand{\headrulewidth}{0pt}
}
\fancypagestyle{intermedio}{
    \lhead{}
    \chead{\fontsize{10}{12}\selectfont\MakeUppercase{\titulo}}
    \rhead{}
    \lfoot{}
    \cfoot{\fontsize{10}{12}\selectfont Página \textbf{\thepage}\ de \textbf{\pageref{LastPage}}}
    \rfoot{}
    \renewcommand{\headrulewidth}{1pt}
}

% Comandos para secciones
\newcommand{\secnumbersection}[1]{
\addtocounter{section}{1}
\section*{CAPÍTULO \thesection \texorpdfstring{\\}\ #1}
\addcontentsline{toc}{section}{CAPÍTULO \thesection : #1}
}
\newcommand{\secnumberlesssection}[1]{
\section*{#1}
\phantomsection
\addcontentsline{toc}{section}{#1}
}

% Nombres
\addto\captionsspanish{\renewcommand{\contentsname}{ÍNDICE DE CONTENIDOS}}
\addto\captionsspanish{\renewcommand{\listfigurename}{ÍNDICE DE FIGURAS}}
\addto\captionsspanish{\renewcommand{\listtablename}{ÍNDICE DE TABLAS}}
\makeatletter
\renewenvironment{thebibliography}[1]
     {\secnumberlesssection{REFERENCIAS BIBLIOGRÁFICAS}
      \@mkboth{\MakeUppercase\bibname}{\MakeUppercase\bibname}%
      \list{\@biblabel{\@arabic\c@enumiv}}%
           {\settowidth\labelwidth{\@biblabel{#1}}%
            \leftmargin\labelwidth
            \advance\leftmargin\labelsep
            \@openbib@code
            \usecounter{enumiv}%
            \let\p@enumiv\@empty
            \renewcommand\theenumiv{\@arabic\c@enumiv}}%
      \sloppy
      \clubpenalty4000
      \@clubpenalty \clubpenalty
      \widowpenalty4000%
      \sfcode`\.\@m}
     {\def\@noitemerr
       {\@latex@warning{Empty `thebibliography' environment}}%
      \endlist}
\makeatother

% Personalizar Tabla de Contenidos
\usepackage{tocloft}
\renewcommand{\cftsecfont}{\fontsize{12}{14}\selectfont\fontspec{\fuente}}
\renewcommand{\cftsubsecfont}{\fontsize{12}{14}\selectfont\fontspec{\fuente}}
\renewcommand{\cftsubsubsecfont}{\fontsize{12}{14}\selectfont\fontspec{\fuente}}

% @@@@@@@@@@@@@@@@@@@@@@@@@@@@@@@@@@@@@@@@@@@@@@@@@@@@@@@@@@@@>
\begin{document}

\pagestyle{portada}
\pagenumbering{roman}
% NOTE: Este archivo contiene la portada, la dedicatoria, los agradecimientos y el resumen.
% __NO ES NECESARIO MODIFICAR ESTE ARCHIVO__, esas se modifican con los comandos que aparecen en main.tex
%@@@@@@@@@@@@@@@@@@@@@@@@@@@@@@@@@@@@@@@@@@@@@@@@@@@@@@@@@@@@@@
\begin{titlepage}
\begin{center}
\noindent
{\fontsize{18}{22}\selectfont UNIVERSIDAD TÉCNICA FEDERICO SANTA MARÍA \\}
{\fontsize{16}{19}\selectfont DEPARTAMENTO DE INFORMÁTICA \\}
{\fontsize{16}{19}\selectfont \MakeUppercase{\ciudad}\ - CHILE \\}
\vspace{1.5cm}
\includegraphics[width=4.41cm,height=3.34cm]{logo/logo.jpg} \\
\vspace{1.5cm}
{\fontsize{20}{24}\selectfont ``\MakeUppercase{\titulo}'' \\}
\vfill
{\fontsize{16}{19}\selectfont \MakeUppercase{\nombrealumno} \\}
\vfill
{\fontsize{16}{19}\selectfont MEMORIA PARA OPTAR AL TÍTULO DE \\}
{\fontsize{16}{19}\selectfont INGENIERO CIVIL EN INFORMÁTICA \\}
\vspace{1.5cm}
{\fontsize{14}{17}\selectfont Profesor Guía: \nombreprofesor \\}
{\fontsize{14}{17}\selectfont Profesor Correferente: \nombrecorreferente \\}
\vspace{2.5cm}
{\fontsize{14}{17}\selectfont \mesexamen\ - \anioexamen \\}
\end{center}
\end{titlepage}

%@@@@@@@@@@@@@@@@@@@@@@@@@@@@@@@@@@@@@@@@@@@@@@@@@@@@@@@@@@@@@@
\newpage
\setcounter{page}{2}
\
\vfill
\vfill
\begin{flushright}
\noindent {\fontsize{16}{19}\selectfont \textbf{DEDICATORIA} \\}
\end{flushright}
\begin{flushright}
\noindent \dedicatoria
\end{flushright}
\vfill
%@@@@@@@@@@@@@@@@@@@@@@@@@@@@@@@@@@@@@@@@@@@@@@@@@@@@@@@@@@@@@@
\newpage
\begin{center}
\noindent {\fontsize{16}{19}\selectfont \textbf{AGRADECIMIENTOS} \\}
\end{center}
\noindent \agradecimientos
\vfill
%@@@@@@@@@@@@@@@@@@@@@@@@@@@@@@@@@@@@@@@@@@@@@@@@@@@@@@@@@@@@@@
\newpage
\secnumberlesssection{RESUMEN}
\vspace{0.3cm}
\noindent \textbf{Resumen---}\resumen \ \\
\vspace{0.3cm} \\
\noindent \textbf{Palabras Clave---}\palabrasclave \ \\
% @@@@@
\vspace{1.2cm} \\
% @@@@@
%\noindent {\fontsize{16}{19}\selectfont \textbf{ABSTRACT}}
%\vspace{1.2cm} \\
\secnumberlesssection{ABSTRACT}
\vspace{0.3cm}
\noindent \textbf{\emph{Abstract}---}\resumeningles \ \\
\vspace{0.3cm} \\
\noindent \textbf{\emph{Keywords}---}\palabrasclaveingles \ \\
%@@@@@@@@@@@@@@@@@@@@@@@@@@@@@@@@@@@@@@@@@@@@@@@@@@@@@@@@@@@@@@


\newpage
\section*{Glosario}

Aquí se deben colocar las siglas mencionadas en el trabajo y su explicación, por orden alfabético. Por ejemplo: \\

{\setlength{\parskip}{0cm} % Para evitar saltar entre cada elemento nombrado.
%Colocar aquí siglas:
DI: Departamento de Informática.

UTFSM: Universidad Técnica Federico Santa María.
}


%Índice de contenidos:
\newpage
\thispagestyle{portada}
\tableofcontents

%Índice de figuras:
\newpage
\thispagestyle{portada}
\listoftables
\listoffigures

\newpage
\pagestyle{intermedio}
\pagenumbering{arabic}
\section*{Introducción}
Debe proporcionar a un lector los antecedentes suficientes para poder contextualizar en general la situación tratada, a través de una descripción breve del área de trabajo y del tema particular abordado, siendo bueno especificar la naturaleza y alcance del problema; así como describir el tipo de propuesta de solución que se realiza, esbozar la metodología a ser empleada e introducir a la estructura del documento mismo de la memoria.

En el fondo, que el lector al leer la Introducción pueda tener una síntesis de cómo fue desarrollada la memoria, a diferencia del Resumen dónde se explicita más qué se hizo, no cómo se hizo.


\newpage
\section{Capítulo \thesection \newline Definición del problema}

\newpage
\section{Capítulo \thesection \texorpdfstring{\\}\ Marco conceptual}

\newpage
\section{Capítulo \thesection \newline Propuesta de solución}

\newpage
\section{Capítulo \thesection \texorpdfstring{\\} Validación de la solución}

\newpage
\section{Capítulo \thesection \texorpdfstring{\\}\ Conclusiones}


\newpage
\secnumberlesssection{ANEXOS}

En los Anexos se incluye todo aquel material complementario que no es parte del contenido de los capítulos de la memoria, pero que permiten a un lector contar con un contenido adjunto relacionado con el tema.


\newpage
% Bibliografía estilo IEEE con las citas "especiales".
\bibliographystyle{IEEEtranSN}
\bibliography{bibliografia}{}

\end{document}
